\par\textbf{Procedimiento completo para comparar tu solución.} Primero separar viscosidad dinámica de cinemática: $\nu=\mu/\rho$. Reynolds usa la cuerda y la velocidad de ensayo; presión dinámica es otra magnitud.

\begin{enumerate}

\item Calcular viscosidad cinemática.

\[\fitmath{\nu=\frac{1.84\times10^{-5}}{1.18}\approx 1.559322\times10^{-5}\,\mathrm{m^{2}/s}}\]

\item Comparar inercia y viscosidad con longitud cuerda.

\[\fitmath{\mathrm{Re}=\frac{28\cdot 0.65}{\nu}}\]

\[\fitmath{\mathrm{Re}=\frac{28\cdot 0.65}{(1.559322034\times10^{-5})}\approx 1167174\,}\]

\item Calcular presión dinámica por separado.

\[\fitmath{q=\frac{1.18\cdot {28}^{2}}{2}\approx 462.56\,\mathrm{Pa}}\]

\end{enumerate}
